\section{第4章 不定积分}
\subsection{不定积分的概念与性质}
\begin{exercise}{424}{原函数关系与二阶求导}
设 $f$ 在 $(1,+\infty)$ 上的一个原函数为 $F(x)=\ln(x+\sqrt{x^2-1})$。求 $f(x)$ 和 $f'(x)$，说明两次求导所用的定义域条件。若把原函数改为 $F(x)+C$，说明所求两个函数是否改变。
\end{exercise}
\begin{exercise}{425}{原函数族与参数退化}
已知 $\ln x$ 是 $f$ 在 $(0,+\infty)$ 上的一个原函数。对 $a>0$，判断 $\ln(ax)$、$a^{-1}\ln x$、$\ln(a+x)$、$\tfrac12(\ln x)^2$ 中哪些仍为 $f$ 的原函数，必要时给出对 $a$ 的限制。用求导与两原函数之差两种方法核对结论。
\end{exercise}
\begin{exercise}{426}{不同表达式的原函数与端点}
证明 $\arcsin(2x-1)$ 与 $2\arcsin\sqrt x$ 在 $(0,1)$ 上为同一函数的原函数，并求两者之差。讨论两个表达式能否连续延拓至 $[0,1]$，以及能否在端点仍按通常意义作为同一函数的可导原函数。
\end{exercise}
\begin{exercise}{427}{分段原函数的拼接条件}
设 $f(x)=3x^2+1$（$x>0$），$f(x)=2x$（$x\le0$）。判断
\[F(x)=\begin{cases}x^3+x+1,&x>0,\\x^2,&x\le0\end{cases}\]
是否为 $f$ 在 $\mathbb R$ 上的原函数，分别检查连续性与导数。将两个分支各自加上任意常数后，是否可能得到全实数范围的原函数？证明结论。
\end{exercise}
\begin{exercise}{428}{原函数存在性与连续性的区别}
设 $F'=f$ 在开区间 $I$ 上成立。判断 $f+F$ 是否必然可导、是否必然连续、是否必然存在原函数、是否必为初等函数，逐项说明依据。证明原函数存在性时，可使用连续函数必有原函数这一结论；不可把 $f$ 的连续性当作已知条件。
\end{exercise}
\begin{exercise}{429}{跳跃间断与原函数拼接}
设 $f(x)=x^2$（$x\ge0$）、$f(x)=\cos x$（$x<0$），另设 $g(x)=x\sin(1/x)$（$x\ne0$）、$g(0)=0$。分别判断两函数在 $\mathbb R$ 上是否存在原函数。对 $f$ 通过两侧原函数的拼接与 $0$ 处导数给出证明；对 $g$ 验证原函数存在定理所需的条件。
\end{exercise}
\begin{exercise}{430}{导函数与原函数的关系辨析}
已知 $\sin x$ 是 $f(x)$ 的一个原函数，求 $\int f(x)\,\dd x$ 与 $f(x)$。再将条件改为 $f'(x)=\sin x$，写出此时 $f$ 及其全部原函数，说明两种条件提供的信息有何区别。
\end{exercise}
\begin{exercise}{431}{切线斜率与初值确定曲线}
一条可导曲线经过 $(\mathrm e,-1)$，且在每个 $x>0$ 处的切线斜率为 $\ln x/x$。求该曲线方程，证明初值条件使积分常数唯一。再求曲线的单调区间与最低点，以检查所得方程与斜率信息相容。
\end{exercise}
\begin{exercise}{432}{不定积分的微分与量的类型}
设 $f$ 在 $\mathbb R$ 上可微。分别求 $\int f'(x)\,\dd x$、$\dd(\int f'(x)\,\dd x)$、$(\int f'(x)\,\dd x)'$，说明三者中哪些是函数族、哪些是微分或确定函数。用 $f(x)=x^2$ 检验符号及积分常数的处理。
\end{exercise}
\begin{exercise}{433}{积分、求导与微分的逆运算}
设 $f\in C^1(I)$，$I$ 为开区间。判断下列等式在通常不定积分记号下是否成立，错误者改正：$\dd(\int f\,\dd x)=f$；$(\int f\,\dd x)'=\int f'\,\dd x$；$\int f'\,\dd x=f$；$\dd(\int f\,\dd x)=f\,\dd x$。说明等式两侧的任意常数为何不能随意遗漏。
\end{exercise}
\begin{exercise}{434}{导数相等与原函数族相等}
设 $f,g$ 在开区间 $I$ 可导，且 $f'=g'$。证明 $f-g$ 为常数，判断是否必有 $f=g$、$f=g+1$、$(\int f\,\dd x)'=(\int g\,\dd x)'$，并判断 $\int\dd f=\int\dd g$ 作为函数族是否成立。若再给出一点 $f(x_0)=g(x_0)$，哪些结论随之加强？
\end{exercise}
\begin{exercise}{435}{嵌套积分微分与常数}
设 $f$ 在开区间可导，求 $\int\dd[\int\dd f(x)]$。逐层说明每次积分引入的常数和微分消去的常数，再判断交换两次任意常数的记号是否会改变最终函数族。
\end{exercise}
\begin{exercise}{436}{复合自变量恢复与定义范围}
设 $f$ 在 $(0,1)$ 可导，且对 $0<t<\pi/2$ 有 $f'(\sin^2t)=\cos2t+\tan^2t$。求 $f(x)$，再由 $f(1/2)=0$ 确定常数。解题时说明代换 $x=\sin^2t$ 在所选区间的值域，不能将 $t$ 与 $x$ 混作同一变量。
\end{exercise}
\begin{exercise}{437}{基本积分公式与定义域分支}
求
\[\int\frac{(x^2-1)\sqrt{1-x^2}-3x}{x\sqrt{1-x^2}}\,\dd x.\]
先化为基本函数之和，再说明最大实定义区间，以及跨过 $x=0$ 时两侧积分常数是否必须相同。
\end{exercise}
\begin{exercise}{438}{三角恒等变形与原函数核验}
求 $\int\cos2x/(\sin^2x\cos^2x)\,\dd x$，先把被积函数拆成两个基本三角函数的平方。说明答案在哪些开区间成立，并对结果求导检验拆分时的符号。
\end{exercise}
\begin{exercise}{439}{半角恒等式与奇点}
求 $\int(1-2\cos x)/(1-\cos2x)\,\dd x$。先写出被积函数的实定义域，再利用 $1-\cos2x=2\sin^2x$ 积分。检查哪些被排除的点为真正奇点，积分常数应如何按区间理解。
\end{exercise}
\begin{exercise}{440}{指数基数与代数约分}
求 $\int[1/(5^x\mathrm e^x)+(1-\mathrm e^{2x})/(1+\mathrm e^x)]\,\dd x$。先说明指数乘积与分式的化简，再确定满足 $F(0)=0$ 的原函数。若将第一项的 $5$ 改为实参数 $a>0$，说明何时需要单独处理常数被积项。
\end{exercise}
\begin{exercise}{441}{基本积分与降幂公式}
求 $\int[6\sqrt x+5\mathrm e^x+\cos^2(x/2)]\,\dd x$，说明余弦平方项的处理。进一步求满足 $F(0)=0$ 的连续原函数，并核验右端点处的右导数。
\end{exercise}
\begin{exercise}{442}{多项式除法与反三角积分}
求 $\int[x^4/(1+x^2)-2/\sqrt{1-x^2}]\,\dd x$。先对有理项作多项式除法，明确共同实定义区间；再确定满足 $F(0)=1$ 的原函数。
\end{exercise}
\begin{exercise}{443}{不连通定义域上的积分常数}
在 $\mathbb R\setminus\{0\}$ 上求 $1/x$ 的全部原函数。证明将答案统一写成 $\ln|x|+C$ 且只允许一个常数会遗漏原函数，并判断是否存在能连续延拓至 $0$ 的原函数。
\end{exercise}
\begin{exercise}{444}{拼接处可导性的必要条件}
设 $f_a(x)=x$（$x<0$）、$f_a(0)=a$、$f_a(x)=2x$（$x>0$）。求使 $f_a$ 在 $\mathbb R$ 上有原函数的全部实数 $a$，并求满足 $F(0)=1$ 的原函数。分别证明条件的必要性和充分性。
\end{exercise}
\subsection{换元积分法}
\begin{exercise}{445}{线性与二次内函数的凑微分}
分别求
\[\int\frac{\dd x}{\sqrt{2x+3}},\qquad \int\frac{x+1}{x^2+2x+5}\,\dd x,\qquad \int\frac{x}{(2x+3)^2}\,\dd x.\]
每式写出代换及实定义区间；比较第三式与第一式为何不能使用相同的幂函数积分步骤直接得到答案。
\end{exercise}
\begin{exercise}{446}{根式内函数与代换值域}
分别求
\[\int\frac{\mathrm e^{2\sqrt x}}{\sqrt x}\,\dd x,\qquad \int\frac{\sin\sqrt x}{\sqrt x}\,\dd x,\qquad \int\frac{\dd x}{(1+x)\sqrt x}.\]
使用同一个根式代换，写出新变量范围。讨论各被积函数在 $x\to0^+$ 时的有限性，以及所得原函数能否连续延拓至 $0$。
\end{exercise}
\begin{exercise}{447}{幂指数、参数与抽象复合换元}
设 $n\in\mathbb R$，$f\in C^1(\mathbb R)$。求
\[\int x^2\mathrm e^{x^3}\,\dd x,\quad \int\frac{x}{1+x^4}\,\dd x,\quad \int\frac{x^2}{(2x^3-1)^n}\,\dd x,\quad \int xf(x^2)f'(x^2)\,\dd x.\]
第三式先限制 $2x^3-1>0$，分别处理 $n=1$ 与 $n\ne1$；第四式说明原函数中 $f$ 应取何种幂次，并以求导验证。
\end{exercise}
\begin{exercise}{448}{三角积分的幂次与降幂选择}
求下列四个不定积分，并分别注明采用的恒等变形或代换：
\[\int\sin^2x\cos^2x\,\dd x,\quad \int\sin^5x\cos^2x\,\dd x,\quad \int\cot^3x\csc x\,\dd x,\quad \int\sin5x\sin3x\,\dd x.\]
对第三式说明为何需要分开理解跨越 $k\pi$ 的积分常数；其余各式给出可在全实数范围成立的原函数。
\end{exercise}
\begin{exercise}{449}{平方和与平方差的参数积分}
对实数 $a\ne0$，求 $\int\dd x/(a^2+x^2)$ 与 $\int\dd x/(a^2-x^2)$，答案须对 $a<0$ 仍成立。说明两式的实定义区间和绝对值的必要性；再单独处理 $a=0$，判断能否直接把它代入非零参数公式。
\end{exercise}
\begin{exercise}{450}{复合指数函数的链式识别}
求 $\int\mathrm e^{\tan x}\sec^2x\,\dd x$ 与 $\int(1-x^{-2})\mathrm e^{x+1/x}\,\dd x$。逐式写出内函数及其导数，注明答案成立的区间，并说明第二式中不能跨越 $x=0$ 统一使用一个积分常数的原因。
\end{exercise}
\begin{exercise}{451}{多层对数代换与实定义域}
分别求
\[\int\frac{\dd x}{x\sqrt{1-\ln^2x}},\quad \int\frac{\dd x}{x\ln x[1+\ln(\ln x)]^2},\quad \int\frac{1+\ln x}{1+x^2\ln^2x}\,\dd x.\]
写出每层代换，并逐式给出实定义域，特别检查第二式中的对数与分母条件。
\end{exercise}
\begin{exercise}{452}{反三角内函数与链式系数}
分别求
\[\int\frac{\dd x}{\sqrt{1-x^2}(\arcsin x)^2},\quad \int\frac{\arcsin\sqrt x}{\sqrt{x(1-x)}}\,\dd x,\quad \int\frac{\arcsin(1-x)}{\sqrt{2x-x^2}}\,\dd x.\]
对每式注明实定义区间；通过求导核验第二式中的系数与第三式中的负号。
\end{exercise}
\begin{exercise}{453}{配方与两个原函数的常数差}
求 $\int\dd x/\sqrt{x-x^2}$，分别使用配方三角代换和 $x=\sin^2t$、$0<t<\pi/2$。说明两种所得表达式在哪一开区间为同一被积函数的原函数，并求两表达式之间的常数差。
\end{exercise}
\begin{exercise}{454}{有界内函数与反正切原函数}
求 $\int\cos x/(9+\sin^2x)\,\dd x$，并确定满足 $F(0)=0$ 的原函数。说明代换 $u=\sin x$ 在全实数轴不单调，为何仍可通过求导确认所得表达式在全实数轴成立。
\end{exercise}
\begin{exercise}{455}{根式分母的拆分代换}
求 $\int(1-x)/\sqrt{9-4x^2}\,\dd x$，分别处理常数分子项与一次分子项。明确实定义区间，求满足 $F(0)=0$ 的原函数，并核验两个部分的系数。
\end{exercise}
\begin{exercise}{456}{三角根式内函数与连通区间}
求 $\int(\sin x+\cos x)/\sqrt[3]{\sin x-\cos x}\,\dd x$，立方根按实数定义。给出被积函数全部实定义区间，并证明原函数在各区间可连续延拓至端点。比较分母改为平方根时实定义域的变化，说明不能给立方根附加非负真数条件。
\end{exercise}
\begin{exercise}{457}{根式三角复合与负号}
求 $\int x\tan(\sqrt{1-x^2})/\sqrt{1-x^2}\,\dd x$。写出代换变量的范围并检验正切是否在此范围产生额外奇点，再讨论被积函数在 $x=\pm1$ 的极限及可补值性。
\end{exercise}
\begin{exercise}{458}{对数根式的复合积分}
求
\[\int\sqrt{\frac{\ln(x+\sqrt{1+x^2})+1}{1+x^2}}\,\dd x.\]
说明对数真数在全实数范围恒为正，再由外层根式确定实定义范围。选择对数整体作新变量，完成积分并确定满足 $F(0)=0$ 的原函数；说明为何不能忽略根式对定义域的额外限制。
\end{exercise}
\begin{exercise}{459}{指数根式代换与补值}
求 $\int\mathrm e^{\sqrt{1+\sin x}}\cos x/\sqrt{1+\sin x}\,\dd x$。在 $1+\sin x>0$ 的各开区间给出原函数，再讨论被排除点处原函数能否连续、可导延拓，说明根式中零点对链式求导的影响。
\end{exercise}
\begin{exercise}{460}{拆项与反正切复合}
求 $\int[x+(\arctan x)^2]/(1+x^2)\,\dd x$，说明为何应先拆分再分别代换。确定满足 $F(0)=1$ 的原函数，并用求导检验平方项积分后应出现的幂次与系数。
\end{exercise}
\begin{exercise}{461}{三角幂与有理化}
求 $\int\sin x\cos^3x/(1+\cos^2x)\,\dd x$，选取代换把它化为有理函数积分。再给出所得原函数的最小正周期，并说明原函数的周期性为何需要实际检查，不能仅由被积函数周期性推出。
\end{exercise}
\begin{exercise}{462}{正切代换的分支与全局原函数}
求 $\int\dd x/(\sin^2x+2\cos^2x)$，先在 $(-\pi/2,\pi/2)$ 使用正切代换。讨论直接使用反正切主值表达式为何在 $\pi/2$ 附近需要调整常数，并说明如何拼接得到全实数范围的一个原函数。
\end{exercise}
\begin{exercise}{463}{指数代换与两种等价表达式}
求 $\int\dd x/(1+\mathrm e^x)$，分别用 $u=\mathrm e^x$ 和 $v=\mathrm e^{-x}$ 代换。证明两种所得原函数只相差常数，并求满足 $F(0)=0$ 的表达式。
\end{exercise}
\begin{exercise}{464}{对数差的整体代换}
求 $\int[\ln(x+1)-\ln x]/[x(x+1)]\,\dd x$。先确定实定义域，计算对数差的导数后整体代换。讨论若改成 $\ln|x+1|-\ln|x|$，原函数公式可延伸到哪些额外区间。
\end{exercise}
\begin{exercise}{465}{对数代换后的多项式拆分}
求 $\int\ln^2x/[x(1+\ln^2x)]\,\dd x$，把代换后的有理函数先拆分再积分。求满足 $F(1)=0$ 的原函数，并证明它在 $(0,+\infty)$ 单调递增，须处理导数为零的点。
\end{exercise}
\begin{exercise}{466}{共轭有理化与根式分支}
求 $\int x/(x-\sqrt{x^2-1})\,\dd x$。先确定实定义域，使用共轭因子化简，分别说明在 $x>1$ 与 $x<-1$ 上计算所得原函数时哪些符号与对数绝对值必须保留。
\end{exercise}
\begin{exercise}{467}{指数复合与二倍角恒等式}
求 $\int\mathrm e^{\sin2x}\sin^2x/\mathrm e^{2x}\,\dd x$。先合并指数，再利用二倍角恒等式将其导数与剩余因子联系，给出在全实数范围成立的原函数，并说明关键系数的来源。
\end{exercise}
\begin{exercise}{468}{对数正切的整体代换}
在 $0<x<\pi/2$ 上求 $\int\ln(\tan x)/\sin2x\,\dd x$。证明 $\ln\tan x$ 在该区间严格递增且值域为 $\mathbb R$，用它作新变量积分，并确定满足 $F(\pi/4)=0$ 的原函数，须核对二倍角产生的系数。
\end{exercise}
\begin{exercise}{469}{商函数求导的逆向识别}
求 $\int[x+\sin x\cos x]/(\cos x-x\sin x)^2\,\dd x$。通过对 $\sin x/(\cos x-x\sin x)$ 求导寻找原函数，并说明答案应限制在哪些开区间；不要求显式求出所有分母零点。
\end{exercise}
\begin{exercise}{470}{三角指数乘积的有理结构}
求 $\int(\cos^2x-\sin x)/\{\cos x[1+\mathrm e^{\sin x}\cos x]\}\,\dd x$。以 $u=\mathrm e^{\sin x}\cos x$ 作整体代换，把积分化为有理函数形式；说明各分母非零条件与对数绝对值的作用，不要求显式求解全部奇点。
\end{exercise}
\begin{exercise}{471}{半角组合与指数乘积导数}
求 $\int(1+\sin x)\mathrm e^x/(1+\cos x)\,\dd x$。用半角恒等式将被积函数写成 $\mathrm e^x$ 与正切半角函数的乘积导数，给出最大实定义区间，并核验原函数在每个区间成立。
\end{exercise}
\begin{exercise}{472}{已知原函数后的复合换元}
在 $0<u<\pi$ 已知 $\int f(u)\,\dd u=\ln(\sin u)+C$。求 $\int xf(1-x^2)\,\dd x$，给出由 $0<1-x^2<\pi$ 确定的实定义区间。解释为何先由已知条件恢复 $f$ 或直接使用原函数换元可得到一致结果。
\end{exercise}
\begin{exercise}{473}{含抽象函数的商法则逆运算}
设 $f\in C^1(I)$，$0\notin I$。证明
\[\int\frac{xf'(x)-(1+x)f(x)}{x^2\mathrm e^x}\,\dd x=\frac{f(x)}x\mathrm e^{-x}+C.\]
明确使用的乘积与商法则，并用 $f(x)=x\mathrm e^x$ 检查两侧。若 $I$ 跨过 $0$，说明公式的适用区间必须如何调整。
\end{exercise}
\begin{exercise}{474}{偶次余项与根式代换}
求 $\int x^3\sqrt{4-x^2}\,\dd x$，使用 $u=4-x^2$ 后先将剩余 $x^2$ 消去。给出实积分区间，并确定满足 $F(0)=0$ 的原函数；说明代换在 $0$ 两侧不单调为何不妨碍最终求导核验。
\end{exercise}
\begin{exercise}{475}{平方和根式与双曲型对数}
求 $\int x^2/(1+x^2)^{3/2}\,\dd x$。分别用三角代换与拆成 $1/\sqrt{1+x^2}-1/(1+x^2)^{3/2}$ 的方法，证明所得答案相差常数；确定满足 $F(0)=0$ 的原函数。
\end{exercise}
\begin{exercise}{476}{平方差根式与参数绝对值}
设 $a>0$，求 $\int x^2/(x^2-a^2)^{3/2}\,\dd x$。在 $x>a$ 上使用三角代换，并给出在 $x<-a$ 上也成立的带绝对值表达式，明确两侧常数可独立选择。
\end{exercise}
\begin{exercise}{477}{根式有理化与变量范围}
求 $\int\dd x/[1+\sqrt[3]{x+1}]$，立方根按实数定义，采用 $t=\sqrt[3]{x+1}$ 并写出 $\dd x$。给出全部最大实定义区间，再检查所得原函数在 $x=-1$ 处是否可导，说明代换的逆函数导数在该点不存在是否妨碍最终答案成立。
\end{exercise}
\begin{exercise}{478}{不同根指数的统一有理化}
求 $\int\dd x/(\sqrt x+\sqrt[3]{x^2})$。先确定实定义区间，再根据两个根指数的最小公倍数选取代换，写出有理化后的完整积分。比较分别只代换平方根或立方根时仍会留下的根式，并检查最终结果的导数。
\end{exercise}
\begin{exercise}{479}{根式分式的有理代换}
求 $\int(3-x)^{-1}\sqrt{(5-x)/(x-1)}\,\dd x$。用根式本身作新变量，写出反解关系及新变量范围；给出被积函数的各个实定义区间，说明 $x=3$ 对积分常数的影响。
\end{exercise}
\begin{exercise}{480}{指数代换后的部分分式}
求 $\int\dd x/(\mathrm e^x+\mathrm e^{2x})$，将它化为新变量的有理函数积分。确定满足 $F(0)=0$ 的原函数，并由被积函数符号判断其单调性。
\end{exercise}
\begin{exercise}{481}{指数加常数的根式代换}
求 $\int\dd x/\sqrt{\mathrm e^x+1}$，令 $t=\sqrt{\mathrm e^x+1}$ 后写出完整有理积分。给出新变量范围和实对数表达式，并说明为何不能把 $\dd x$ 误写为 $2t\dd t$。
\end{exercise}
\begin{exercise}{482}{幂代换与不连通区间}
求 $\int\dd x/[x(x^6+1)]$。分别使用 $t=x^6$ 与先拆成 $1/x-x^5/(x^6+1)$ 的方法，比较所得表达式。解释 $x>0$、$x<0$ 两个区间中 $\ln|x|$ 和积分常数的处理。
\end{exercise}
\begin{exercise}{483}{倒代换与根式绝对值}
求 $\int\dd x/[x^4\sqrt{x^2+1}]$。尝试 $t=1/x$，明确 $\sqrt{x^2+1}$ 变换后出现的绝对值，分别在正负半轴积分；最终给出可在两个半轴分别使用的表达式并求导核验。
\end{exercise}
\begin{exercise}{484}{对数与代数的混合内函数}
求 $\int(1+x)/(x^2+x\ln x)\,\dd x$。选择 $u=x+\ln x$ 化简，证明分母在正半轴只有一个零点，从而说明答案成立的最大开区间，不要求写出该零点的初等表达式。
\end{exercise}
\begin{exercise}{485}{配方、三角代换与半圆结构}
求 $\int\sqrt{2x-x^2}\,\dd x$，先配方并给出实定义范围。确定满足 $F(1)=0$ 的原函数，并利用被积函数关于 $x=1$ 的对称性检验所得原函数的奇对称关系。
\end{exercise}
\begin{exercise}{486}{指数根式与两步代换}
求 $\int\dd x/\sqrt{1+\mathrm e^{2x}}$，先作指数代换，再选用根式或三角代换。写出对数形式的实数原函数，说明所有对数真数在全实数范围的符号并求导核验。
\end{exercise}
\begin{exercise}{487}{三角代换后的有理积分}
求 $\int\dd x/[(1+x^2)\sqrt{1-x^2}]$。在 $(-1,1)$ 作 $x=\sin t$，再使用正切代换完成积分。明确两次代换的变量范围，避免反函数分支导致常数跳跃。
\end{exercise}
\begin{exercise}{488}{参数退化与复合换元}
对 $a\in\mathbb R$ 求 $\int x/(1+ax^2)\,\dd x$，完整处理 $a=0$。在包含 $0$ 的最大实定义区间上选取 $F_a(0)=0$ 的原函数，并求 $a\to0$ 时 $F_a(x)$ 的点态极限，检验参数退化的一致性。
\end{exercise}
\begin{exercise}{489}{反三角主值与分支常数}
在 $(-1,1)$ 上求 $\int\dd x/\sqrt{1-x^2}$。若代换 $x=\sin t$ 分别取 $t\in(-\pi/2,\pi/2)$ 与 $t\in(\pi/2,3\pi/2)$，写出两种计算过程，并说明 $\sqrt{\cos^2t}=|\cos t|$ 与反函数分支如何共同保证所得函数只相差常数。
\end{exercise}
\subsection{分部积分法}
\begin{exercise}{490}{多项式与指数三角函数的分部积分}
求 $\int x\mathrm e^{ax}\,\dd x$、$\int x^2\cos x\,\dd x$、$\int(x^2-x)\sin x\,\dd x$，其中 $a\in\mathbb R$。第一式须单独处理 $a=0$，其余两式明确每次分部积分中次数降低的部分，并通过求导核验符号。
\end{exercise}
\begin{exercise}{491}{反函数与对数函数的分部选择}
求 $\int x/\sin^2x\,\dd x$、$\int\ln(x+\sqrt{x^2+1})\,\dd x$、$\int\arctan x\,\dd x$。写出每式分部积分中 $u,\dd v$ 的选择，说明第一式的定义区间与第二式对数真数的符号。
\end{exercise}
\begin{exercise}{492}{循环积分、恒等式与抽象函数}
求 $\int\mathrm e^{-x}\cos x\,\dd x$ 与 $\int x\cot^2x\,\dd x$。另设 $\int f(x)\,\dd x=(1+\sin x)\ln x+C$（$x>0$），求 $\int xf'(x)\,\dd x$；说明第三式为何可利用已知原函数而不必完整展开 $f'$。
\end{exercise}
\begin{exercise}{493}{对数移位与多项式分部}
求 $\int x\ln(x+1)\,\dd x$，给出实定义区间。分别采用直接分部积分与先令 $u=x+1$ 的方法，验证两种答案之差为常数，再确定满足 $F(0)=0$ 的原函数。
\end{exercise}
\begin{exercise}{494}{平移反正弦与根式符号}
求 $\int\arcsin(x+1)\,\dd x$，明确可导原函数的开区间。将答案连续延拓至实定义域端点，求两个端点的单侧导数，说明被积函数在端点有限与自身导数无界的区别。
\end{exercise}
\begin{exercise}{495}{凑微分与分部积分的配合}
求 $\int x\cot x/\sin^2x\,\dd x$，先把 $\cot x/\sin^2x$ 识别为一个基本函数的导数，再进行分部积分。注明积分成立的区间，并解释只作 $u=\cot x$ 代换为何仍会留下需要处理的 $x$。
\end{exercise}
\begin{exercise}{496}{半角恒等式后的分部积分}
求 $\int x\cos^4(x/2)/\sin^3x\,\dd x$。先用半角公式化简三角因子，再选择分部积分；明确原被积函数的定义域，并检查化简是否在原来无定义的点造成表面约分。
\end{exercise}
\begin{exercise}{497}{对数三角函数与分部积分}
在 $0<x<\pi$ 上求 $\int\ln(\sin x)/\sin^2x\,\dd x$。通过分部积分把剩余项化为基本三角积分，并用求导检验。说明若要延伸到 $\sin x<0$ 的区间，题面中的对数必须作何调整。
\end{exercise}
\begin{exercise}{498}{根式积分的两种路线}
求 $\int x^3/\sqrt{1+x^2}\,\dd x$，分别用整体根式代换与取 $u=x^2$、$\dd v=x\dd x/\sqrt{1+x^2}$ 的分部积分法。证明两种结果等价，并说明本题中哪一步使剩余积分的次数降低。
\end{exercise}
\begin{exercise}{499}{指数三角循环积分}
用两次分部积分求 $\int\mathrm e^{2x}\sin3x\,\dd x$。进一步推导实参数 $a,b$ 满足 $a^2+b^2>0$ 时 $\int\mathrm e^{ax}\sin bx\,\dd x$ 的公式，并检查 $b=0$ 与 $a=0$ 的退化情形。
\end{exercise}
\begin{exercise}{500}{平方和根式与参数退化}
设 $a>0$，用分部积分求 $\int\sqrt{x^2+a^2}\,\dd x$，将出现的同类积分移项求解，说明对数真数的正性。再单独求 $a=0$ 时的 $\int|x|\,\dd x$，并讨论两个半轴上的表达式如何拼接。
\end{exercise}
\begin{exercise}{501}{平方差根式与区间约束}
设 $a>0$，求 $\int\sqrt{a^2-x^2}\,\dd x$，分别用三角代换与分部积分核验。明确实定义域和反正弦的主值范围，再确定满足 $F(0)=0$ 的原函数及其奇偶性。
\end{exercise}
\begin{exercise}{502}{反正弦与剩余根式积分}
求 $\int x\arcsin x\,\dd x$，明确分部积分后的根式积分如何继续化简。选取满足 $F(0)=0$ 的原函数，检验它的奇偶性，并求两端连续延拓值。
\end{exercise}
\begin{exercise}{503}{反正切乘负幂与对数绝对值}
求 $\int x^{-2}\arctan x\,\dd x$，说明 $\int\dd x/[x(1+x^2)]$ 的分解。分别给出正负半轴上的原函数族，并讨论被积函数在 $x\to0$ 的主项是否允许原函数有限连续延拓。
\end{exercise}
\begin{exercise}{504}{对数分部与内部极点}
求 $\int\ln(1+x)/(2-x)^2\,\dd x$。先给出最大实定义区间，再作分部积分并处理所得有理积分。明确 $\ln(1+x)$ 与 $\ln|2-x|$ 的适用条件，说明 $x=2$ 两侧的常数为何独立。
\end{exercise}
\begin{exercise}{505}{指数根式的分部与有理化}
在 $x>0$ 上求 $\int x\mathrm e^x/\sqrt{\mathrm e^x-1}\,\dd x$。先把指数根式凑为一个微分，再对剩余根式积分换元。确定能连续延拓至 $0$ 且取值为 $0$ 的原函数，并检查右导数。
\end{exercise}
\begin{exercise}{506}{对数与根式的双重处理}
求 $\int\ln x/\sqrt{1+x}\,\dd x$（$x>0$），用分部积分后作 $t=\sqrt{1+x}$。说明对数中 $t-1$ 的符号，并讨论所求原函数在 $x\to0^+$ 时是否具有有限极限。
\end{exercise}
\begin{exercise}{507}{反正切根式与端点补值}
求 $\int\arctan\sqrt x\,\dd x$，给出在 $x>0$ 上的计算过程。选择满足 $F(0)=0$ 的连续延拓，核验其在 $0$ 的右导数与被积函数在该点的值相等。
\end{exercise}
\begin{exercise}{508}{有理指数乘积的导数识别}
求 $\int x^2\mathrm e^x/(x+2)^2\,\dd x$。把被积函数写成 $\mathrm e^x$ 与一个有理函数的导数结构，可设原函数为 $\mathrm e^x(A+B/(x+2))$ 求系数。注明定义区间，并用求导检查该设定确实覆盖本题。
\end{exercise}
\begin{exercise}{509}{三角有理组合与逆向求导}
求 $\int[(1+x)\sin x+\cos x]/(1+x)^2\,\dd x$。分别用分部积分的消项与识别商函数导数的方法，说明两条路线的等价性，并明确 $x=-1$ 两侧的积分常数。
\end{exercise}
\begin{exercise}{510}{复合导函数恢复与分部积分}
设 $f'(\mathrm e^t)=\cos t$ 对所有实数 $t$ 成立。求 $f(x)$ 的全部表达式，并由 $f(1)=0$ 确定常数。解题时先写出 $f'$ 的定义区间，再处理含 $\cos\ln x$ 的积分。
\end{exercise}
\begin{exercise}{511}{正割幂的递推与初值}
在不含 $\cos x=0$ 的开区间上，令 $I_n=\int\sec^nx\,\dd x$。对整数 $n\ge2$ 证明
\[I_n=\frac{\sin x}{(n-1)\cos^{n-1}x}+\frac{n-2}{n-1}I_{n-2},\]
其中等式按相差任意常数理解。求 $I_0,I_1$ 并用递推式求 $I_3,I_4$，说明递推为何需要两个初始积分。
\end{exercise}
\begin{exercise}{512}{参数分部积分与零参数极限}
对实参数 $a$ 求 $\int x\mathrm e^{ax}\,\dd x$，选取满足 $F_a(0)=0$ 的原函数。单独处理 $a=0$，并用泰勒公式证明非零参数表达式在 $a\to0$ 时趋于该结果，解释积分常数选择对极限的影响。
\end{exercise}
\begin{exercise}{513}{对数幂积分的递推}
设 $n$ 为非负整数，$J_n=\int(\ln x)^n\,\dd x$（$x>0$）。建立 $J_n$ 与 $J_{n-1}$ 的递推关系，并求 $J_3$。证明所得表达式可取有限的 $x\to0^+$ 极限，不能仅以每个对数项都无界判断原函数发散。
\end{exercise}
\subsection{有理函数的积分}
\begin{exercise}{514}{假分式除法与余式积分}
求 $\int x^3/(x^2+1)\,\dd x$，先写出多项式除法的商与余式。确定满足 $F(0)=0$ 的原函数，再用奇偶性检验该原函数的对称性。
\end{exercise}
\begin{exercise}{515}{不同一次因子的部分分式}
求 $\int(3x-2)/(x^2+3x-4)\,\dd x$，写出分母分解和部分分式系数。给出所有最大实定义区间，说明所得对数的绝对值及各区间常数的含义。
\end{exercise}
\begin{exercise}{516}{三不同实根分母的积分}
求 $\int(2x^2-x+1)/[(x+1)(x-3)(x+2)]\,\dd x$。用赋值法或待定系数法完整确定三个一次分式系数，再以通分核验分子恒等式，明确原函数的各个定义区间。
\end{exercise}
\begin{exercise}{517}{简单极点与原函数族}
求 $\int\dd x/[x(x-1)]$。在三个最大定义区间上分别描述全部原函数，判断是否能找到可连续延拓至 $0$ 或 $1$ 的原函数，并由对数主项说明理由。
\end{exercise}
\begin{exercise}{518}{重一次因子的完整分解}
求 $\int(x^2+1)/[(x+1)(x^2-1)]\,\dd x$。先按重根结构列全部分分式，说明若只保留 $1/(x+1)$ 与 $1/(x-1)$ 两项为何不足，再求系数并积分。
\end{exercise}
\begin{exercise}{519}{不可约二次因子的分子形式}
求 $\int\dd x/[(x^2+1)(1+2x)]$。写出包含二次因子对应一次分子的完整分解，积分后同时保留对数与反正切项。说明若把该一次分子擅自设为常数，会遗漏什么约束。
\end{exercise}
\begin{exercise}{520}{配方与分子拆成导数}
求 $\int(3x+1)/(x^2+2x+5)\,\dd x$，把分子拆成分母导数的倍数与常数。证明分母对所有实数非零，并确定满足 $F(-1)=0$ 的原函数。
\end{exercise}
\begin{exercise}{521}{重二次分母与标准积分}
求 $\int\dd x/(x^2+2x+5)^2$，先配方再使用三角代换或分部积分。写出有理项和反正切项，最后对整个表达式求导核验平方分母的系数。
\end{exercise}
\begin{exercise}{522}{四次分母分解与对称分子}
求 $\int(2x^3+2x^2+5x+5)/(x^4+5x^2+4)\,\dd x$。先分解分母，再利用分子的结构减少待定系数数量；说明为何答案可在全实数轴使用同一个积分常数。
\end{exercise}
\begin{exercise}{523}{配方后重二次因子的积分}
求 $\int x^2/(x^2+2x+2)^2\,\dd x$。先作平移代换，把分子拆成适合积分的三项，并给出完整原函数。通过检查大 $|x|$ 时被积函数的衰减阶，判断所求原函数在两端是否具有有限极限。
\end{exercise}
\begin{exercise}{524}{高次有理函数的实因式分解}
求 $\int(x^4+1)/(1+x^6)\,\dd x$。利用 $1+x^6=(x^2+1)(x^4-x^2+1)$，将四次因子进一步分解为两个实不可约二次因子，再作部分分式积分。注明实定义区间，并利用被积函数的偶性检查可取为奇函数的原函数。
\end{exercise}
\begin{exercise}{525}{高重零点与偶函数分解}
求 $\int\dd x/(x^4-x^6)$。先利用偶性寻找仅含 $1/x^2$、$1/x^4$ 与 $1/(1-x^2)$ 的分解，再积分。列出全部被排除点，说明不能将 $x=0$ 当作普通的简单极点处理。
\end{exercise}
\begin{exercise}{526}{倒幂代换与重因子结构}
求 $\int(1+2x^4)/[x^3(1+x^4)^2]\,\dd x$。先作 $t=1/x^2$ 以降低次数，再积分得到实原函数，分别说明正负半轴上的适用性，并通过求导核验重因子对应的有理项。
\end{exercise}
\begin{exercise}{527}{四次平方和的实分解}
求 $\int\dd x/(x^4+1)$，将分母分解为两个实不可约二次因子后积分。给出在 $0$ 附近连续且取值为 $0$ 的表达式，并讨论合并两个反正切为一个主值反正切时可能需要补加的常数。
\end{exercise}
\begin{exercise}{528}{原函数存在性与奇偶性}
逐项判断并证明或举例：初等函数在其连续的定义区间内存在原函数；不连续函数也可能存在原函数；奇函数在关于 $0$ 对称的开区间上的原函数都是偶函数；偶函数在这样的区间上的原函数都是奇函数。对最后一项给出使之成立的积分常数条件，并说明“存在原函数”与“有初等原函数”有何区别。
\end{exercise}
\begin{exercise}{529}{闭区间连续性保证的结论}
设 $f\in C[a,b]$，$a<b$。判断 $f$ 在 $(a,b)$ 内是否必有导函数、原函数、最大值或最小值、局部极值。对必然成立者说明定理依据；对其余结论给出具体反例，注意闭区间最值定理的取值范围。
\end{exercise}
\begin{exercise}{530}{分段函数原函数的参数拼接}
设 $f(x)=2(x-1)$（$x<1$）、$f(x)=\ln x$（$x\ge1$）。求实数 $a,b$ 使
\[F_{a,b}(x)=\begin{cases}(x-1)^2,&x<1,\\x(\ln x+a)+b,&x\ge1\end{cases}\]
为 $f$ 在 $\mathbb R$ 上的原函数。必须分别检查拼接处连续、左右导数及非拼接处导数，不能仅由两侧分别可积确定结论。
\end{exercise}
\begin{exercise}{531}{反正切指数函数的分部积分}
求 $\int\mathrm e^{-2x}\arctan(\mathrm e^x)\,\dd x$。先作指数代换，再通过分部积分处理反正切项，明确新变量恒正对对数绝对值的影响。
\end{exercise}
\begin{exercise}{532}{反正切与有理因子的配合}
求 $\int\arctan x/[x^2(1+x^2)]\,\dd x$。先分解 $1/[x^2(1+x^2)]$，再分别采用分部与整体代换；写出在两个半轴上的原函数族，并讨论 $0$ 处是否可连续延拓。
\end{exercise}
\begin{exercise}{533}{指数有理混合的对数导数}
求 $\int(1+x)/[x(1+x\mathrm e^x)]\,\dd x$，尝试将被积函数写成两个对数导数之差。证明 $1+x\mathrm e^x$ 在全实数范围恒正，从而给出答案的全部最大实定义区间。
\end{exercise}
\begin{exercise}{534}{对数商函数的逆向求导}
求 $\int(1-\ln x)/(x-\ln x)^2\,\dd x$。尝试对 $x/(x-\ln x)$ 求导识别原函数，并用最值证明分母在正半轴不为零，确定全实定义区间。
\end{exercise}
\begin{exercise}{535}{指数根式与反正切分部积分}
在 $x>0$ 上求 $\int\mathrm e^{2x}\arctan\sqrt{\mathrm e^x-1}\,\dd x$。先令 $t=\sqrt{\mathrm e^x-1}$ 化为多项式乘反正切的积分，再作分部积分。明确代换后 $\mathrm e^{2x}\dd x$ 的表达，避免遗漏雅可比因子。
\end{exercise}
\begin{exercise}{536}{三角商函数的导数结构}
求 $\int[x+\sin x\cos x]/(x\sin x+\cos x)^2\,\dd x$。通过求导 $\cos x/(x\sin x+\cos x)$ 寻找原函数，说明答案的适用开区间，不要求用初等函数写出分母的所有零点。
\end{exercise}
\begin{exercise}{537}{三角有理函数与万能代换}
求 $\int\dd x/(\sin2x+2\sin x)$。先因式分解并确定全部奇点，再作 $t=\tan(x/2)$ 积分。说明万能代换遗漏的点是否都是原函数的奇点，以及积分常数如何在定义区间上选择。
\end{exercise}
\begin{exercise}{538}{万能代换与可去的参数断点}
求 $\int\dd x/(1+\sin x+\cos x)$，用 $t=\tan(x/2)$ 化为有理函数。列出原被积函数的全部奇点，区别原函数的奇点与代换坐标在奇数倍 $\pi$ 处的失效，说明如何检验答案。
\end{exercise}
\begin{exercise}{539}{反正切原函数的跨区间拼接}
求 $\int\dd x/(2+\cos x)$，先在 $(-\pi,\pi)$ 采用万能代换。原被积函数在全实数连续，说明为何得到的反正切主值公式仍需按周期补常数，并给出拼接为全实数原函数的方法。
\end{exercise}
\begin{exercise}{540}{常数项与对数导数拆分}
求 $\int(3\cos x-\sin x)/(\cos x+\sin x)\,\dd x$，把分子拆成分母及其导数的线性组合。给出每个最大定义区间的原函数，并明确对数中绝对值的必要性。
\end{exercise}
\begin{exercise}{541}{对称拆分的三角积分}
求 $\int\sin x/(\sin x+\cos x)\,\dd x$，利用分母与其导数拆分。再求对应的 $\int\cos x/(\sin x+\cos x)\,\dd x$，检验两者相加与被积函数恒等式一致，说明常数的合并方式。
\end{exercise}
\begin{exercise}{542}{根式与线性项的三角代换}
求 $\int\dd x/[2x+\sqrt{1-x^2}]$，在 $(-1,1)$ 使用 $x=\sin t$。求出分母在该区间的唯一零点，分别给出两侧原函数，并说明代换中根式为何可写为 $\cos t$。
\end{exercise}
\begin{exercise}{543}{奇偶拆分与重二次分母}
求 $\int(x^3+1)/(1+x^2)^2\,\dd x$，先把奇函数部分与偶函数部分分开处理。确定满足 $F(0)=0$ 的原函数，并说明最终答案中哪些项分别来自这两部分。
\end{exercise}
\begin{exercise}{544}{反正切指数与根式分母}
求 $\int x\mathrm e^{\arctan x}/(1+x^2)^{3/2}\,\dd x$，先令 $t=\arctan x$，注意 $\cos t$ 的符号，再完成指数三角积分。说明该原函数可在全实数范围成立的依据。
\end{exercise}
\begin{exercise}{545}{指数几何和的因式分解}
求 $\int\dd x/(1+\mathrm e^{x/2}+\mathrm e^{x/3}+\mathrm e^{x/6})$。根据指数分母的最小公倍数选取 $t=\mathrm e^{x/6}$ 化为有理积分，利用 $1+t+t^2+t^3=(1+t)(1+t^2)$ 作部分分式分解，再积分并给出实定义范围，须明确 $\dd x$ 中的系数。
\end{exercise}
\begin{exercise}{546}{三角因子与有理代换}
求 $\int\dd x/[(2+\cos x)\sin x]$，先确定全部实定义区间，再用 $u=\cos x$ 化为有理函数积分。解释代换时 $\sin^2x=1-u^2$ 的作用，并核验所得对数系数。
\end{exercise}
\begin{exercise}{547}{反正弦分部与根式多项式}
求 $\int x^3\arcsin x/\sqrt{1-x^2}\,\dd x$。先求不含反正弦因子的原函数，再作分部积分，明确实定义区间。说明若先令 $t=\arcsin x$，为何会转化为多项式三角幂的分部积分。
\end{exercise}
\begin{exercise}{548}{隐函数约束下的原函数}
设 $y=y(x)$ 是关系 $y(x-y)^2=x$ 在某开区间上的可导分支，且 $x-3y\ne0$。求 $\int\dd x/(x-3y(x))$，可引入 $u=x-y$ 作参数。推导 $x,y$ 关于 $u$ 的表达式与 $\dd x$，并明确参数化需要排除的值及积分公式的分支适用条件。
\end{exercise}
\begin{exercise}{549}{函数关系的单射区间与换元}
设 $f(\sin^2t)=t/\sin t$ 对 $0<t<\pi/2$ 成立。求 $\int\sqrt x\,f(x)/\sqrt{1-x}\,\dd x$，明确 $f$ 在 $(0,1)$ 上由条件唯一确定的依据。先作 $x=\sin^2t$，再完成分部积分，并用直接恢复 $f(x)$ 的方法核验。
\end{exercise}
\begin{exercise}{550}{已知原函数与高阶导数积分}
设 $\int f(x)\,\dd x=\mathrm e^{x^2}+C$。求 $f(x)$、$f'(x)$ 及 $\int xf''(x)\,\dd x$。最后一个积分先用分部积分写为 $xf'(x)-f(x)$ 再代入，说明为什么不必展开 $f''$。
\end{exercise}
\begin{exercise}{551}{加权原函数恢复与奇点}
已知 $\int xf(x)\,\dd x=\arcsin x+C$ 在 $(-1,0)$ 或 $(0,1)$ 上成立。求 $f(x)$ 及 $\int f(x)\,\dd x$，给出带绝对值的实表达式。说明在包含 $0$ 的开区间上是否可能存在连续 $f$ 满足原条件。
\end{exercise}
\begin{exercise}{552}{指数复合导函数的恢复}
设 $\int f'(\mathrm e^t)\,\dd t=-(1+t)\mathrm e^{-t}+C$，且 $f(1)=0$。求 $f(x)$（$x>0$），分别说明对已知积分式求导和将 $x=\mathrm e^t$ 代回的步骤，再用初值确定常数。
\end{exercise}
\begin{exercise}{553}{抽象函数商的平方导数}
设 $f\in C^2(I)$，且 $f'(x)\ne0$。证明
\[\int\left[\frac{f(x)}{f'(x)}-\frac{f^2(x)f''(x)}{[f'(x)]^3}\right]\,\dd x=\frac12\left[\frac{f(x)}{f'(x)}\right]^2+C.\]
逐项检查商法则与链式法则产生的分母幂次，并解释条件 $f'\ne0$ 的必要性；再用 $f(x)=\mathrm e^x$ 验证。
\end{exercise}
\begin{exercise}{554}{线性三角组合的平方分母}
设实数 $a,b$ 满足 $a^2+b^2>0$。证明
\[\int\frac{\dd x}{(a\sin x+b\cos x)^2}=-\frac1{a^2+b^2}\frac{a\cos x-b\sin x}{a\sin x+b\cos x}+C.\]
给出结果成立的区间条件，并分别检查 $a=0$、$b=0$。再说明若 $a=b=0$，为何问题本身失去定义。
\end{exercise}
\begin{exercise}{555}{参数因子与部分分式退化}
对 $a\in\mathbb R$ 求 $\int\dd x/[(x-1)(x-a)]$，完整区分 $a=1$ 与 $a\ne1$。再取 $0<a<2$，在公共区间 $x<0$ 选取满足 $F_a(-1)=0$ 的原函数，求 $a\to1$ 时 $F_a(x)$ 的极限，说明积分常数归一化对该极限的作用。
\end{exercise}
\begin{exercise}{556}{一次因子与重二次因子的联合分解}
求 $\int\dd x/[x(x^2+1)^2]$，先列全部分分式并利用奇偶性减少系数。给出两个半轴上的全部原函数，最后通过求导及 $x\to0$ 的主项检查对数系数是否正确。
\end{exercise}
